\section*{Capítulo \ref{ch:g7:earthquakes} --- Os terremotos}

\begin{solution}{exo:g7:earthquakes:1}
(a) \omterm{def:g7:earthquakes:focus}{foco}; (b) \omterm{def:g7:earthquakes:intensity}{intensidade}; (c) \omterm{def:g7:earthquakes:fault}{falha}; (d) \omterm{def:g7:earthquakes:focus}{epicentro}; (e) \omterm{def:g7:earthquakes:magnitude}{magnitude};
(f) \omterm{def:g7:earthquakes:seismometer}{sismômetro}.
\end{solution}

\begin{solution}{exo:g7:earthquakes:2}
(a) Na própria aldeia: o \omterm{def:g7:earthquakes:focus}{epicentro} é o ponto da superfície bem em cima do
\omterm{def:g7:earthquakes:focus}{foco}. (b) \qty{30}{km}, bem para baixo. (c) Em geral menor: o tremor
enfraquece com a distância. (Ela também depende do terreno: sobre \omterm{def:g3:stones-pebbles-sand:sand}{lama}
mole, a segunda aldeia ainda poderia tremer forte.)
\end{solution}

\begin{solution}{exo:g7:earthquakes:3}
(a) \omterm{def:g7:earthquakes:magnitude}{Magnitude} (um número com uma casa decimal). (b) \omterm{def:g7:earthquakes:intensity}{Intensidade} (o que
foi visto num lugar). (c) \omterm{def:g7:earthquakes:intensity}{Intensidade} (algarismos romanos, diferentes em
dois lugares). (d) \omterm{def:g7:earthquakes:magnitude}{Magnitude}: a energia liberada, um degrau de \omterm{def:g7:earthquakes:magnitude}{magnitude}
a mais.
\end{solution}

\begin{solution}{exo:g7:earthquakes:4}
(a) Falso: os cientistas não conseguem prever os \omterm{def:g4:volcanoes-earthquakes-first-look:earthquake}{terremotos}; eles só
sabem que \omterm{def:g7:earthquakes:fault}{falhas} vão se romper de novo um dia. (b) Falso: em pleno
\omterm{def:g2:water-on-earth:ocean}{oceano}, um \omterm{def:g7:earthquakes:tsunami}{tsunami} raramente passa de um metro de altura; ele cresce
perto do \omterm{def:g1:landscapes-around-us:water}{litoral}. (c) Verdadeiro. (d) Verdadeiro: 92~\% dos \omterm{def:g4:volcanoes-earthquakes-first-look:earthquake}{terremotos} de
\omterm{def:g7:earthquakes:magnitude}{magnitude}~6 ou mais acontecem a menos de \qty{300}{km} de um \omterm{def:g7:plate-tectonics:plate-boundary}{limite de
placas}.
\end{solution}

\begin{solution}{exo:g7:earthquakes:5}
(a) Cerca de 127 por ano de \omterm{def:g7:earthquakes:magnitude}{magnitude} 6 a 6.9; cerca de 13 de \omterm{def:g7:earthquakes:magnitude}{magnitude} 7
a 7.9. (b) $126.7 \div 13.12 \approx 9.7$: cerca de 10~vezes menos. (c)
$13.12 + 0.76 + 0.04 = 13.92$: cerca de \textbf{14} por ano. (d)
$0.04 \times 50 = 2$: dois em 50~anos (Sumatra 2004 e Japão 2011).
\end{solution}

\begin{solution}{exo:g7:earthquakes:6}
(a) Um degrau: 32~vezes. (b) Dois degraus: $32 \times 32 = 1024$, cerca
de 1000~vezes. (c) Cada degrau deixa as ondas 10~vezes maiores: $10 \times 10 =
100$~vezes maiores.
\end{solution}

\begin{solution}{exo:g7:earthquakes:7}
(a) $8 \times 5 = \qty{40}{km}$. (b) $8 \times 25 = \qty{200}{km}$. (c)
$120 \div 8 = \qty{15}{s}$. (d) Uma estação dá só uma distância, não uma
direção: todos os pontos a essa distância formam um círculo em volta da
estação, e o \omterm{def:g7:earthquakes:focus}{epicentro} poderia estar em qualquer lugar dele.
\end{solution}

\begin{solution}{exo:g7:earthquakes:8}
(a) V. (b) II. (c) VIII. (d) A \omterm{def:g7:earthquakes:intensity}{intensidade} muda de um lugar para outro: o
relato (c) vem de um lugar perto do \omterm{def:g7:earthquakes:focus}{epicentro}, ou construído sobre
terreno mole; o relato (a), de um lugar mais distante, ou sobre \omterm{def:g3:stones-pebbles-sand:rock}{rocha}
sólida. O \omterm{def:g4:volcanoes-earthquakes-first-look:earthquake}{terremoto} continua tendo uma única \omterm{def:g7:earthquakes:magnitude}{magnitude}.
\end{solution}

\begin{solution}{exo:g7:earthquakes:9}
(a) $7031 - 5561 = 1470$. (b) $1470 \div 7031 \approx 0.209$: cerca de
\textbf{21~\%}. (c) A \omterm{def:g7:plate-tectonics:plate}{placa} que afunda desce bem fundo no \omterm{def:g6:inside-the-earth:layers}{manto}. Ela
continua rígida o bastante para se romper, então \omterm{def:g4:volcanoes-earthquakes-first-look:earthquake}{terremotos} podem começar
dentro dela, a centenas de quilômetros de profundidade.
\end{solution}

\begin{solution}{exo:g7:earthquakes:10}
(a) $2500 \div 800 = 3.125$~horas: 3~horas e $0.125 \times 60 = 7.5$
minutos, cerca de 3~horas e 8~minutos. (b) 3~h~8~min $-$ 2~h~17~min $=$
51~minutos depois. (c) A onda levou horas para chegar, enquanto um grande
\omterm{def:g4:volcanoes-earthquakes-first-look:earthquake}{terremoto} embaixo do \omterm{def:g2:water-on-earth:ocean}{mar} é localizado em poucos minutos: um alerta
poderia ter mandado as pessoas para lugares altos muito antes de a onda
chegar.
\end{solution}

\begin{solution}{exo:g7:earthquakes:11}
(a) $400 \div 800 = 0.5$~hora: 30~minutos. (b) $20 \div 40 = 0.5$~hora:
30~minutos. (c) 60~minutos. (d) Os últimos \qty{20}{km} levam tanto tempo
quanto os primeiros \qty{400}{km}: a onda desacelera enormemente nas
águas rasas. (e) A frente da onda desacelera enquanto a água atrás dela
ainda vem rápido: a água se amontoa, e a onda fica mais alta.
\end{solution}

\begin{solution}{exo:g7:earthquakes:12}
(a) Três degraus: $32 \times 32 \times 32 = 32\,768$, cerca de
\textbf{33\,000}~vezes. (b) $32\,768 \div 127 \approx 258$: cerca de
260~anos de todos os \omterm{def:g4:volcanoes-earthquakes-first-look:earthquake}{terremotos} de \omterm{def:g7:earthquakes:magnitude}{magnitude}~6 do mundo. (c) Quase toda a
energia é liberada pelos poucos maiores \omterm{def:g4:volcanoes-earthquakes-first-look:earthquake}{terremotos}: uma única \omterm{def:g7:earthquakes:magnitude}{magnitude}~9
pesa mais que séculos de \omterm{def:g4:volcanoes-earthquakes-first-look:earthquake}{terremotos} de \omterm{def:g7:earthquakes:magnitude}{magnitude}~6.
\end{solution}

\begin{solution}{pb:g7:earthquakes:1}
\textbf{1.} A: $45 - 25 = 20$~s. B: $29 - 16 = 13$~s. C: $34 - 19 =
15$~s.

\textbf{2.} A estação B: o seu atraso é o mais curto, e quanto mais curto
o atraso, mais perto o \omterm{def:g4:volcanoes-earthquakes-first-look:earthquake}{terremoto}. (As suas ondas rápidas também chegaram
primeiro.)

\textbf{3.} A: $8 \times 20 = \qty{160}{km}$. B: $8 \times 13 =
\qty{104}{km}$. C: $8 \times 15 = \qty{120}{km}$.

\textbf{4.} Todos os pontos a \qty{160}{km} de A: um círculo de raio
\qty{160}{km} com o centro em A. O \omterm{def:g7:earthquakes:focus}{epicentro} fica em algum lugar dele.

\textbf{5.} A: $160 \div 20 = 8$~quadrados. B: $104 \div 20 = 5.2$~quadrados.
C: $120 \div 20 = 6$~quadrados.

\textbf{6.} E fica $3 + 6 = 9$~quadrados a leste e 14.5~quadrados ao norte
de O, no \omterm{def:g2:water-on-earth:ocean}{mar}, entre as três estações; os círculos de A e de B passam por
ele.

\textbf{7.} Embaixo do \omterm{def:g2:water-on-earth:ocean}{mar}. Um movimento repentino do fundo do \omterm{def:g2:water-on-earth:ocean}{mar} pode
provocar um \omterm{def:g7:earthquakes:tsunami}{tsunami}: os \omterm{def:g1:landscapes-around-us:water}{litorais} e a ilha precisam ser alertados.

\textbf{8.} Um degrau de \omterm{def:g7:earthquakes:magnitude}{magnitude}: cerca de 32~vezes mais energia.

\textbf{9.} $12 \times 20 = \qty{240}{km}$.

\textbf{10.} $240 \div 800 = 0.3$~hora, e $0.3 \times 60 = 18$~minutos.

\textbf{11.} 5~minutos depois do \omterm{def:g4:volcanoes-earthquakes-first-look:earthquake}{terremoto}.

\textbf{12.} $18 - 5 = 13$: os moradores de Port Hale têm cerca de
\textbf{13~minutos} para chegar a um lugar alto. (A onda desacelera nas
águas rasas perto da ilha, o que dá um pouco mais de tempo, mas ninguém
deveria contar com isso: eles partem imediatamente.)
\end{solution}
